\section{Introduction}\label{sec:intro}

A wireless sensor network (WSN) consists of battery-powered nodes that sense their surroundings and forward readings to a sink, usually over several hops. Every packet a node sends or relays costs energy, and once the first node is exhausted part of the network can be cut off. The number of rounds until the first node dies (FND) is therefore a standard measure of network lifetime~\cite{chang2004maximum}, and lifetime-aware routing is a central design problem in WSNs and in the low-power Internet of Things (IoT).

A growing line of work applies deep reinforcement learning (DRL) to this problem: neural networks learn to choose next hops, cluster heads or transmit powers~\cite{yang2025madii,younus2021sdwsn,tang2025twotier}. These methods need long training and often a GPU. They are usually evaluated against clustering heuristics, metaheuristics or other learned methods. We found no recent learned WSN router that was compared with a battery-weighted shortest-path router, and none that reported its lifetime as a fraction of a provable optimum. Recent comparison studies even tabulate published results obtained on different datasets and in different units~\cite{icmcsi2026}. As a result, it is hard to tell how much the learning contributes. Neighbouring areas of networking show the same pattern: simple, well-tuned baselines match deep models once they are evaluated carefully~\cite{simplemodel2025}.

This paper answers the question by measurement. The maximum-lifetime linear program (LP) of Chang and Tassiulas~\cite{chang2004maximum} gives a ceiling that no routing scheme can exceed. For traffic that varies over time we use an oracle version that knows all future traffic. Every method is scored as the percentage of that ceiling it reaches, on held-out deployments and on real traffic. Against this yardstick the methods separate into clear levels. A recent multi-agent DRL router, MADII~\cite{yang2025madii}, rebuilt in its own published setting, reaches about 11--15\% of the ceiling. A published energy-and-load router, EBR-DA~\cite{mahdi2018ebrda}, reaches about 15\%. Battery-weighted Dijkstra reaches 85--88\% with no training.

The remaining gap comes from coordination: in battery-weighted Dijkstra every node picks its path independently, so in the same round many nodes can pick the same relay. We propose \emph{LEST load-table Dijkstra}, a solver-free router that closes about half of the gap. Each round the sink builds the routing tree one node at a time, nearest the sink first. Each node's path is found by Dijkstra over a link cost that combines radio energy, the exact residual battery and the load already booked on each relay in that round. Each source's traffic rate reaches the sink through a Lightweight State Table (LEST): every node reports its load as one of four tiers, piggybacked on packets it already sends. Finally, an LP planner re-solved every round from live batteries serves as an upper reference for what full flow optimisation adds.

\paragraph{Contributions} The contributions of this paper are:
\begin{enumerate}[label=(C\arabic*), leftmargin=*]
  \item \textbf{An oracle-normalised benchmark for lifetime routing.} Every method is reported as a percentage of a provable ceiling: the maximum-lifetime LP for static traffic, and an oracle LP, solved by bisection, for time-varying traffic. The benchmark covers 30 held-out deployments per scenario, four traffic scenarios, real event-driven traffic from the Intel Berkeley Lab dataset~\cite{intellab2004} and paired Wilcoxon tests with Holm correction.
  \item \textbf{A measured ranking of learned, published and classical routers.} MADII, rebuilt in its own setting and checked against its published figures, reaches 11.1\% of the ceiling (19.5 rounds against the paper's 19). Battery-weighted Dijkstra reaches 88.4\%, and DRL warm-started from Dijkstra adds no significant gain ($p=0.24$).
  \item \textbf{LEST load-table Dijkstra.} A per-round, solver-free and training-free router that builds the tree sequentially with within-round load booking, using exact energy and a 4-tier shared load table whose overhead is charged in the energy model. It reaches 90.5\% of the oracle against 85.0\% for battery-weighted Dijkstra (118/1/1 wins, ties and losses over 120 deployments) and 90.8\% against 86.0\% on Intel Lab traffic.
  \item \textbf{Ablations and negative results.} Both the energy term and the load term are needed (load only 45.1\%, energy only 85.0\%, both 90.5\%), and energy must stay exact rather than tiered. Five ways of adding prediction, lifetime weighting or LP prices to Dijkstra stay within about $\pm1$ point of it. A Holt-Winters forecast adds only 0.5 points to an LP planner that is re-solved every round.
\end{enumerate}
All code, the result files and the scripts that generate every table and figure are released (see the data availability statement).

Section~\ref{sec:related} reviews related work and positions this paper. Section~\ref{sec:model} states the system model, the packet-layer model and every assumption. Section~\ref{sec:methods} describes the proposed router and the reference methods, and Section~\ref{sec:setup} the evaluation protocol. Section~\ref{sec:results} reports the results and the statistical tests, Section~\ref{sec:discussion} discusses them together with the threats to validity, and Section~\ref{sec:conclusion} concludes.
